\documentclass[UTF8,a4paper,11pt,openany]{ctexbook}
\usepackage{../common/statlearnbook}
\SLSetBookMetadata{第 1 册}{统计学习方法讲义 v2.6.0}{数学基础与统计学习基本理论}
\begin{document}
\setcounter{page}{184}
\setcounter{chapter}{10}
\pagestyle{headings}
\chaptermark{模型评估、选择与泛化}
\section*{交叉验证}
简单留出法把数据分为训练、验证与测试三部分；\kFold{}交叉验证把用于开发的数据划为$K$个互不相交折，每轮用一个折验证、其余折训练，最后平均验证损失；当$K=N$时称留一交叉验证。分类问题通常使用分层划分维持类别比例，分组或时间相关数据则必须按组或时间切分，不能假设样本可任意打乱。

学习算法不仅包含单次模型拟合，还包含预处理、特征选择、超参数搜索和决策阈值选择。所有从数据估计的步骤都必须放入训练折内部，否则验证折的信息会提前进入模型。

% v2.3.1 figure UID: UFIG-P158-01
% Proposed post-v2.3.1 polish for UFIG-P158-01: protect key-label size and stroke hierarchy.
\tikzset{slfig-UFIG-P158-01/.style={font=\fontsize{9.2pt}{11.0pt}\selectfont,
  every path/.append style={line cap=round,line join=round}}}
% Unnumbered, leakage-free model-selection workflow. The two-row layout avoids global scaling.
\begin{center}
\begin{tikzpicture}[slfig-UFIG-P158-01,
  flowstep/.style={slfig node,line width=.76pt,minimum width=18mm,text width=16mm,
    minimum height=10.5mm,inner xsep=2.5mm,inner ysep=1.6mm,
    font=\fontsize{9.2pt}{11pt}\selectfont},
  mini/.style={slfig node accent,line width=.72pt,minimum width=14mm,text width=12mm,
    minimum height=9mm,inner xsep=2.2mm,inner ysep=1.5mm,
    font=\fontsize{9.2pt}{11pt}\selectfont},
  badge/.style={circle,fill=SLBlue,text=white,minimum size=5.2mm,inner sep=0pt,
    font=\bfseries\fontsize{9.2pt}{11pt}\selectfont},
  mainflow/.style={-{Stealth[length=2.2mm,width=1.4mm]},draw=SLBlue,
    line width=.98pt,shorten >=.65mm},
  feedback/.style={-{Stealth[length=1.9mm,width=1.2mm]},draw=SLGray,
    line width=.72pt,dash pattern=on 3.2pt off 2.2pt,shorten >=.8mm},
  testflow/.style={-{Stealth[length=2.1mm,width=1.35mm]},draw=SLRed,
    line width=.88pt,shorten >=.75mm},
]
  % First row: freeze the rule, create development folds, and keep every fit inside its fold.
  \node[flowstep] (raw) at (0,0) {原始数据};
  \node[badge] at ([xshift=-1mm,yshift=1mm]raw.north west) {1};
  \node[flowstep] (split) at (24mm,0) {固定切分};
  \node[badge] at ([xshift=-1mm,yshift=1mm]split.north west) {2};
  \node[mini] (prep) at (47mm,0) {预处理};
  \node[mini] (select) at (63mm,0) {特征选择};
  \node[mini] (fit) at (79mm,0) {拟合};
  \node[mini] (valid) at (95mm,0) {折内验证};
  \node[badge] at ([xshift=-1mm,yshift=1mm]prep.north west) {3};
  \node[flowstep] (choose) at (117mm,0) {汇总验证损失\\选择候选};
  \node[badge] at ([xshift=-1mm,yshift=1mm]choose.north west) {4};

  \draw[mainflow] (raw)--(split);
  \draw[mainflow] (split)--(prep);
  \draw[mainflow] (prep)--(select);
  \draw[mainflow] (select)--(fit);
  \draw[mainflow] (fit)--(valid);
  \draw[mainflow] (valid)--(choose);

  \node[slfig compact note,line width=.56pt,minimum width=62mm,minimum height=6mm,
    font=\fontsize{9.2pt}{11pt}\selectfont]
    (folds) at (71mm,-12.5mm)
    {$K$ 折：每一折都重新执行完整流水线，不复用折外拟合状态};
  \draw[feedback] (choose.south) |- (36mm,-18.5mm) |- (prep.west);
  \node[font=\fontsize{9.2pt}{11pt}\selectfont,text=SLGray]
    at (71mm,-22.5mm) {反馈仅留在开发区};

  % Second row is read right-to-left: retrain on all development data, then unlock one test use.
  \node[flowstep] (refit) at (88mm,-34mm) {全部开发集\\重新训练};
  \node[badge] at ([xshift=-1mm,yshift=1mm]refit.north west) {5};
  \node[flowstep] (final) at (20mm,-34mm) {最终一次评估};
  \node[badge] at ([xshift=-1mm,yshift=1mm]final.north west) {6};
  \draw[mainflow] (choose.south) |- (refit.east);
  \draw[mainflow] (refit.west)--(final.east);

  \node[slfig locked,line width=.78pt,rounded corners=1.7mm,minimum width=26mm,minimum height=11mm,
        align=center,font=\fontsize{9.2pt}{11pt}\selectfont] (test) at (-5mm,-34mm)
        {\textbf{锁定测试集}\\仅最终一次使用};
  \draw[mainflow] (split.south) |- (test.north);
  \draw[testflow]
    (test.east)--(final.west);

  \begin{scope}[on background layer]
  \node[draw=SLTeal,fill=SLTealBG,line width=.72pt,rounded corners=2mm,
        fit=(prep)(select)(fit)(valid)(folds)(choose)(refit),inner sep=3.5mm,
        label={[font=\bfseries\fontsize{9.2pt}{11pt}\selectfont,text=SLTeal]
          above:开发集（允许循环选择）}] (dev) {};
  \node[draw=SLRed,fill=SLRedBG,line width=.86pt,
        dash pattern=on 3.2pt off 2.2pt,rounded corners=2mm,
        fit=(test),inner sep=2.5mm,
        label={[font=\bfseries\fontsize{9.2pt}{11pt}\selectfont,text=SLRed]
          below:测试区（禁止反馈）}] {};
  \end{scope}
\end{tikzpicture}
\end{center}


\paragraph{首次阅读检查：无泄漏的模型选择。}先按输入与状态、初始化、一次更新、保持量或目标、停止条件与输出证书读取原文。任何由数据估计的变换都只能在当前训练折内拟合；测试集只在候选锁定后使用一次。
\end{document}
